\chapter*{CAPÍTULO 3. SocialScrum: Proceso ágil para la gestión de la deuda social}
\addcontentsline{toc}{chapter}{CAPÍTULO 3. SocialScrum: Proceso ágil para la gestión de la deuda social}
\setcounter{chapter}{3}
\setcounter{section}{0}
\setcounter{subsection}{0}
\setcounter{subsubsection}{0}
\renewcommand{\thesection}{\thechapter.\arabic{section}}
\renewcommand{\thesubsection}{\thesection.\arabic{subsection}}
\renewcommand{\thesubsubsection}{\thesubsection.\arabic{subsubsection}}

Este capítulo para su mejor comprensión se encuentra dividido en tres secciones principales. La primera sección introducimos el termino ``SocialScrum'' y la necesidad de su desarrollo en los equipos de desarrollo de software y una visión general del proceso ágil a realizar. En la segunda sección, describimos el modelo propuesto directamente incluyendo sus componentes clave. En la tercera sección, pasamos a detallar los fundamentos teóricos y sus valores que le dan origen y la justificación de su desarrollo, por último, en la cuarta sección, detallamos el diseño del modelo SocialScrum, describiendo sus componentes clave y cómo se integran para abordar la deuda social en equipos de desarrollo de software.

\section{Introducción a SocialScrum}
SocialScrum es un proceso ágil que adapta los principios, eventos y artefactos de Scrum para identificar, priorizar y mitigar la deuda social en equipos de desarrollo. Su propósito es hacer visibles las dinámicas socio-técnicas que erosionan la colaboración y la calidad, integrando prácticas de inspección y adaptación en el plano humano y organizacional [CITA NECESARIA]. En síntesis, promueve un entorno de trabajo saludable y sostenible donde las mejoras sociales y del producto evolucionan de forma armónica. 

\section{Modelo SocialScrum}
El modelo SocialScrum se compone de cuatro componentes principales:
\begin{itemize}
    \item \textbf{Eventos:} Los eventos de SocialScrum son los mismos que los de Scrum, pero adaptados para abordar la deuda social.
    \item \textbf{Artefactos:} Los artefactos de SocialScrum son los mismos que los de Scrum, pero adaptados para abordar la deuda social.
    \item \textbf{Roles:} Los roles de SocialScrum son los mismos que los de Scrum, pero adaptados para abordar la deuda social.
    \item \textbf{Prácticas:} Las prácticas de SocialScrum son las mismas que las de Scrum, pero adaptadas para abordar la deuda social.
\end{itemize}

\section{Fundamentos teóricos y valores}

\section{Diseño del modelo SocialScrum}

\section{Conclusiones}